%! Function Model Selection and Construction — Unit 1, Lesson 1.7 — Homework (Answer Key)
\documentclass[10pt]{article}
\usepackage{precalculus-article}
\usepackage{precalculus-key}   % pulls in -boxes; do NOT also load -boxes
\newcommand{\TallMath}[1]{$\displaystyle #1\rule[-1.4em]{0pt}{3.2em}$}

\begin{document}

\pageheader{Unit 1, Lesson 1.7}{Homework --- Answer Key}

\vspace{0.12in}

\begin{enumerate}[itemsep=10pt]

\item Take differences until they are constant, then name each model type.

\smallskip
\renewcommand{\arraystretch}{1.4}
\begin{center}
\begin{tabular}{|c|c|c|c|c|c|}
\hline
$x$ & 0 & 1 & 2 & 3 & 4 \\
\hline
$y$ & 11 & 8 & 5 & 2 & $-1$ \\
\hline
1st & --- & \ans{$-3$} & \ans{$-3$} & \ans{$-3$} & \ans{$-3$} \\
\hline
\end{tabular}
\qquad
\begin{tabular}{|c|c|c|c|c|c|}
\hline
$x$ & 0 & 1 & 2 & 3 & 4 \\
\hline
$y$ & 1 & 3 & 11 & 31 & 69 \\
\hline
1st & --- & \ans{2} & \ans{8} & \ans{20} & \ans{38} \\
\hline
2nd & --- & --- & \ans{6} & \ans{12} & \ans{18} \\
\hline
3rd & --- & --- & --- & \ans{6} & \ans{6} \\
\hline
\end{tabular}
\end{center}

\smallskip
Left table: \ans{linear} \qquad Right table: \ans{cubic}

\item Match each context to its model type. Write a letter in each blank:
      \textbf{L} linear, \textbf{Q} quadratic, \textbf{C} cubic,
      \textbf{R} rational, \textbf{P} piecewise.

\begin{enumerate}[label=(\alph*), itemsep=4pt]
  \item A pool is filled at a steady 12 gallons per minute. \ans{L}
  \item The amount of paint needed to cover a square wall, as a function of the
        wall's side length. \ans{Q}
  \item A city charges \$0.05 per gallon of water for the first 5{,}000 gallons
        and \$0.09 per gallon after that. \ans{P}
  \item The time to drive 300 miles, as a function of average speed
        (time $=$ 300/speed). \ans{R}
  \item The amount of concrete needed to fill a cube-shaped footing, as a
        function of the footing's edge length. \ans{C}
\end{enumerate}

\item \textbf{A rectangle with no wall to lean on.} A gardener has 100 feet of
      fencing for a rectangular plot and must fence \textbf{all four sides}.

\begin{work}
   2w + 2\ell &= 100 \\
         \ell &= 50 - w \\
         A(w) &= w(50 - w) \\
              &= 50w - w^2 \\
            w &= -\frac{50}{2(-1)} \\
              &= 25
\end{work}

\begin{enumerate}[label=(\alph*), itemsep=4pt]
  \item Constraint: \ans{$2w + 2\ell = 100$} \qquad so $\ell = $ \ans{$50 - w$}
  \item $A(w) = $ \ans{$50w - w^2$} \qquad Domain: \ans{0} $< w <$ \ans{50}
  \item The largest area occurs at $w = $ \ans{25} ft, which makes the plot
        a \ans{square}. Explain in one sentence why this answer is a square but
        the barn pen from class was not.

        \ansline{Here all four sides cost fence, so width and length are treated alike; at the barn one side was free, so it paid to make that side longer.}
\end{enumerate}

\boxguard[24]
\item \textbf{Back to the pen.} The garden club's model is
      $A(w) = 80w - 2w^2$ on $0 < w < 40$.

\begin{work}
   A(5) &= 80(5) - 2(5)^2 \\
        &= 400 - 50 \\
        &= 350 \\
  A(35) &= 80(35) - 2(35)^2 \\
        &= 2800 - 2450 \\
        &= 350 \\
  \frac{A(15) - A(5)}{15 - 5} &= \frac{750 - 350}{10} \\
                              &= 40
\end{work}

\begin{enumerate}[label=(\alph*), itemsep=4pt]
  \item $A(5) = $ \ans{350} ft$^2$ \qquad $A(35) = $ \ans{350} ft$^2$.
        Two different widths, the same area --- why does that make sense?

        \ans{Both are 15 ft from the best width $w = 20$, and a parabola is symmetric about its vertex.}
  \item Average rate of change from $w = 5$ to $w = 15$: \ans{40},
        with units \ans{square feet per foot of width}.
  \item Is $w = 45$ in the domain of this model? \ans{No} \quad Why not?
        \ans{Two widths of 45 ft already need 90 ft of fence.}
\end{enumerate}

\item \textbf{An inverse-proportion model.} The force $F$ between two magnets is
      inversely proportional to the square of the distance $d$ between them, so
      $F = \dfrac{k}{d^2}$. At $d = 3$ cm the force measures 20 newtons.

\begin{work}
    20 &= \frac{k}{3^2} \\
     k &= 180 \\
  F(6) &= \frac{180}{6^2} \\
       &= 5
\end{work}

\begin{enumerate}[label=(\alph*), itemsep=4pt]
  \item $k = $ \ans{180} \qquad so $F(d) = $ \ans{$180/d^2$}
  \item $F(6) = $ \ans{5} newtons
  \item Domain restriction: $d$ \ans{$> 0$}, because \ans{a distance cannot be zero or negative}
\end{enumerate}

\item A table of equally spaced inputs has second differences that are all
      equal to 4. Which statement is correct?

\begin{enumerate}[label=\Alph*., itemsep=3pt]
  \item Linear, because the differences are constant.
  \item Quadratic, because the second differences are constant.
        \quad \textcolor{keyred}{\textbf{$\leftarrow$ correct}}
  \item Degree 5, because the table has five rows.
  \item No polynomial model fits, because the second differences are not zero.
\end{enumerate}

\end{enumerate}

\vspace{0.1in}

\boxguard
\begin{extensionbox}
The barn wall turns out to be only \textbf{30 feet long}. If the pen's length
$\ell$ is more than 30 feet, the part beyond the barn has to be fenced too ---
so the fencing equation changes.

\begin{enumerate}[label=(\alph*), itemsep=4pt]
  \item When $\ell > 30$, the fence must cover $2w + \ell + (\ell - 30) = 80$.
        Solve that for $\ell$: \ans{$\ell = 55 - w$}
  \item That case applies when $w$ is \textbf{less than} \ans{25} feet.
  \item Write the two area rules. \quad $A(w) = w(55 - w)$ for
        \ans{$0 < w < 25$}, \quad and $A(w) = w(80 - 2w)$ for \ans{$25 \le w < 40$}.
  \item Both rules give the same area at the threshold. Check it:
        \ans{750} ft$^2$. Why must a piecewise model agree at the threshold?

        \ansline{One width can only have one area --- if the rules disagreed at $w = 25$, the model would not be a function there.}
\end{enumerate}
\end{extensionbox}

\vspace{0.1in}

\boxguard
\begin{notesbox}{Preview of Lesson 2.1}
Every model you built today came out a polynomial, and the useful one came out
quadratic. Unit 2 opens with \textit{Quadratic Functions Revisited}, and
$A(w) = 80w - 2w^2$ is still on the board: where the vertex formula
$-\frac{b}{2a}$ comes from, what vertex form shows that standard form hides, and
why a parabola is symmetric about its vertex --- which is exactly why $A(5)$ and
$A(35)$ came out equal in problem 4.
\end{notesbox}

\end{document}
